% ---------------------------------------------------------------------
% Limitations subsection — immunology contribution to the Discussion.
% ~350 words. Honest disclosure of populations where REAL/RareShield
% cannot make reliable detection claims at Kang-atlas scale.
% Author: Immunology (agent-mozus3vv). Cross-references Math's
% predicted-detectability table (agent/mathematics@f564cc6).
% Target: sections/08_discussion.tex, as a \subsection after the
% "Implications for immunology" subsection (immuno_discussion.tex).
% ---------------------------------------------------------------------

\subsection{Boundaries of what our framework can currently guarantee}
\label{sec:discussion-limitations-immunology}

Three concrete limitations shape the claims we make. First, the Le~Cam
detectability envelope (\Cref{sec:minimax-detectability}) places a
sharp boundary on which rare populations can be recovered label-free
at any given atlas scale. At the Kang pan-cancer immune atlas
($n = 4.9 \times 10^6$, $B = 30$, per-batch heterogeneity
$\Lambda = 3$), populations with prevalence-separation product
$\pi\Delta < 10^{-3}\sigma$ fall below the detection floor, producing
a true-positive rate under 0.3 at $\alpha = 0.05$. Specifically,
\textit{CXCR5}$^+$\textit{TCF7}$^+$ progenitor-exhausted CD8 T
(TPEX), LAMP3$^+$ mregDC, KIR$^+$ adaptive NK, endothelial tip
cells, and drug-tolerant-persister tumour cells all sit in this
regime at Kang alone. For these populations our framework does not
claim detection from the Kang atlas in isolation; their recovery
requires either auxiliary protein or TCR signal, or pooling across
multiple atlases beyond $10^8$ cells. We disclose this explicitly
and present the populations instead as below-floor targets whose
detection-regime analysis (Extended Data~Fig.~1) itself is a
contribution.

Second, some rare populations are lost before integration rather than
during integration. Neutrophil subsets including LOX-1$^+$
PMN-MDSC and TAN-1 through TAN-5 subsets are systematically
undersampled by the standard 10x Genomics droplet-capture protocol,
and mast cells are similarly depleted relative to their tissue
abundance. Our framework is a post-integration detector and does not
and cannot recover populations that never entered the cell-by-gene
matrix. We flag this as a pre-integration identifiability problem
orthogonal to overcorrection, to be addressed by protocol choice
and by modality-aware capture strategies (FACS-sort enrichment prior
to library preparation; split-pool-barcoded whole-transcriptome
methods that capture neutrophils without selection bias).

Third, our evaluation uses annotated rare populations with their
labels hidden at test time. This is a legitimate form of ground
truth but it is not a direct test on unannotated populations.
Novel populations detected under our framework in the pan-cancer
atlas --- for example, a candidate CXCR5-dim TPEX sub-state
\citep{ourfig5} --- require the full evidence package of
\Cref{sec:evidence-package} plus independent cohort replication
before a definite claim of previously-unannotated biology. Main-text
detections are reported only when $\geq 4$ of the six evidence
criteria are satisfied.

\paragraph*{Telescoping detection recovers below-whole-atlas populations at sub-compartment scale.}
The $\pi\cdot\Delta^2 < 1.5\times10^{-4}$ below-floor rule above is
stated at the whole-atlas level and applies to flat REAL. Hierarchical
REAL, applied at lineage-restricted sub-compartments, moves the
operating point from whole-atlas prevalence to compartment-level
prevalence; at each level the flat floor still applies, but the
$\pi\cdot\Delta^2$ product is typically $10$--$100\times$ larger
because rare populations are often concentrated in specific lineages.
This recovers pulmonary ionocytes at the airway-epithelial sub-compartment
level, and the CXCR5-dim TPEX sub-state at the exhausted-CD8 sub-compartment
level (Supplementary Theorem~S1-hierarchical, Mathematics agent
\citealp{math2026hierarchical}). Our disclosure is therefore
three-tier: populations recovered by flat whole-atlas REAL (above
floor at $\pi\cdot\Delta^2 > 1.5\times10^{-4}$, e.g., MAIT, $\gamma\delta$
T, TREM2$^+$ LAM, myCAF, EMT-hybrid), populations recovered by
hierarchical sub-compartment REAL (above floor at sub-compartment
product, e.g., ionocytes, TPEX sub-state, AS-DC), and populations
unresolved at all compartment depths (endothelial tip cells, drug-tolerant
persister baseline state). Main-text claims are restricted to the
first two tiers; tier-three populations are disclosed as present in
the dataset but beyond the current framework's detection capability.

% ---------------------------------------------------------------------
% Requires: sec:minimax-detectability, sec:evidence-package,
% \citep{ourfig5} is a placeholder for the Fig 5 companion section /
% preprint. Biological Science may prefer \Cref{fig:fig5} or
% \Cref{sec:results-tpex-substate} depending on skeleton.
% ---------------------------------------------------------------------
